
d'où le résultat d'%
après (6.1.2) en passant au quotient.

Si @@M0@@, @@M1@@, donc @@M2@@ est un homomorphisme de @@M3@@-algèbres~; de plus il est compatible aux augmentations.

Il reste enfin à voir que c'est un @@M4@@-homomorphisme. Il faut donc montrer la commutativité du diagramme :
@@M5@@

Par additivité, il suffit de la montrer pour les éléments de @@M6@@, pour lesquels elle est évidente, et pour les éléments d'augmentation nulle de @@M7@@, pour lesquels elle résulte, d'après la définition de @@M8@@ et (6.1.4), de celle du diagramme (3.7.2).

\textbf{Proposition 6.9.} \uline{Soit} @@M9@@ \uline{une} @@M10@@-@@M11@@-\uline{algèbre augmentée}, @@M12@@~; \uline{alors} :
@@M13@@
\uline{i}.\uline{e}.
@@M14@@

Il faut voir que ces deux éléments donnent même image dans @@M15@@, c'est-à-dire que la @@M16@@-ième classe de Chern de @@M17@@ est l'image de @@M18@@.

On peut écrire @@M19@@ comme une somme @@M20@@, avec @@M21@@, @@M22@@ pour tout @@M23@@, et @@M24@@. On a alors :
@@M25@@
Comme @@M26@@ est un @@M27@@-homomorphisme, on peut écrire :
@@M28@@
